\documentclass[11pt]{article}
\usepackage[margin=2.3cm]{geometry}
\usepackage{amsmath,amssymb,booktabs}
\usepackage[hidelinks]{hyperref}
\usepackage{microtype}
\usepackage{url}
\DeclareUrlCommand\lean{\urlstyle{tt}\def\UrlBreaks{\do\.\do\_\do\/}}
% the counts of a committee
\newcommand{\NT}{N_{\mathcal T}}
\newcommand{\NP}{N_{\mathrm{pair}}}
\newcommand{\Th}{\mathrm{Th}}
\newcommand{\Fn}[1]{\mathcal F_{\mathrm{#1}}}
\emergencystretch=3em
\setlength{\parskip}{0.35em}
\title{The claims of Threshold leanXMSS}
\author{}
\date{November 2026}
\begin{document}
\maketitle

\noindent This document states, in plain mathematics, every end claim about Threshold leanXMSS in
the repository. Each claim is a statement. This document does not say which claims have a proof. The
claims are these.
\begin{itemize}\itemsep0pt
\item The key generation (the keygen) of the $n = 2t - 1$ operators securely computes the trusted
dealer of PRAWNS with selective abort.
\item Each tree ceremony securely computes the dealer of its tree with selective abort. That dealer
gives the trees of leanXMSS when the circuits compute the tweakable hash.
\item The MPC gadgets of Goyal and Song (coins, multiplications, the check and one batch) each
securely compute their functionality. Each honest run gives the output of its functionality.
\item The threshold PRF is correct and adaptively secure.
\item The two-round signing is correct. It is strongly unforgeable under adaptive corruption in the
random-oracle model.
\item A last claim composes the real keygen, the real ceremonies with retries and the signing
service over the lifetime of the key. It is stated against a static coalition, up to the advantage of
one fixed forger.
\end{itemize}

Each claim is a Lean~4 proposition. The proposition \lean{ThresholdLeanXmssSecurityStatement} in
\lean{Statement.lean} bundles the eighteen component claims. The theorem
\lean{thresholdLeanXmss_security} in \lean{ThresholdLeanXmssSecurity.lean} has the bundle as its
type. Its \texttt{\#guard\_msgs} axiom guard is the test of the proof. The component claims count as
proved only when the axiom guard of \lean{thresholdLeanXmss_security} lists only \lean{propext},
\lean{Classical.choice} and \lean{Quot.sound}, with no \lean{sorryAx} and no added axiom. A guard
that lists \lean{sorryAx} certifies no proof. The lifetime claim \lean{LifetimeSecurityStatement} is a
separate proposition and is not in the bundle. The entry file states no theorem of it, so the guard of
\lean{thresholdLeanXmss_security} does not test it. A proof of the lifetime claim needs its own
theorem and its own guard.

\subsection*{Where to read it}

The statement layer is the folder \lean{formal/ThresholdLeanXmss/}. It imports Mathlib, VCVio and
the leanXMSS statement of leanVM (\lean{formal/lean_xmss} at the pin of \lean{lakefile.toml}), and
nothing of the proofs.
\begin{enumerate}\itemsep0pt
\item \lean{Scheme.lean} holds the protocol as programs of one operator. These are the keygen of
\lean{doc/dkg.md}, the gadgets, the ceremony, the threshold PRF and the signing of
\lean{doc/sign.md}. The end of the file holds the oracles of the world and the execution model.
\item \lean{Ideal.lean} holds the functionalities, the trusted dealer \lean{Fkeygen}, the tree
dealer \lean{Ftree}, and the oracles of the PRAWNS games and of the lifetime game.
\item \lean{Reductions.lean} holds the explicit simulators \lean{keygenSimulator} and
\lean{treeSimulator}, the reductions \lean{prgReduction} and \lean{lifetimeRunReduction} to the AES
term, and the fixed forger \lean{lifetimeForger}. A bound that holds the advantage of a reduction is
only as strong as that reduction, so a review reads this file too.
\item \lean{Statement.lean} holds the counted worlds, the notions of secure computation, the games,
the error terms and the claims. Its module doc has the claims table, the model and the deviations
from the papers. It is the file that a human reviews.
\end{enumerate}
The entry file is \lean{formal/ThresholdLeanXmssSecurity.lean}. Its theorem
\lean{thresholdLeanXmss_security} has the type \lean{ThresholdLeanXmssSecurityStatement}. That type
quantifies over every committee, every check field with its constants, every pair of primitives
$(\Th, \mathrm{AES})$, every circuit family and every setting.
\begin{center}\footnotesize
\begin{tabular}{lll}
\toprule
Claim & In words & Section\\\midrule
\lean{KeygenCompletenessStatement} & an honest keygen gives the keys of $\Fn{keygen}$ & \ref{sec:keygen}\\
\lean{KeygenSecurityStatement} & the keygen securely computes $\Fn{keygen}$ & \ref{sec:keygen}\\
\lean{TreeCompletenessStatement} & an honest ceremony gives the tree of $\Fn{tree}$ & \ref{sec:tree}\\
\lean{TreeSecurityStatement} & each ceremony securely computes $\Fn{tree}$ & \ref{sec:tree}\\
\lean{DealerTreeStatement} & $\Fn{tree}$ gives the trees of $\mathsf{KeyGen}$ & \ref{sec:tree}\\
\lean{TreeSuccessStatement} & $\Fn{tree}$ fails only on a failed search & \ref{sec:tree}\\
\lean{CoinStatement}, \lean{CoinCompletenessStatement} & $\mathsf{Coin}$ and $\Fn{coin}$ & \ref{sec:gadgets}\\
\lean{MultStatement} & $\mathsf{MulRep}$ and $\mathsf{Reshare}$ and $\Fn{mult}$ & \ref{sec:gadgets}\\
\lean{ExtendMultStatement} & $\mathsf{MulSh}$ and $\Fn{extendMult}$ & \ref{sec:gadgets}\\
\lean{MultVerifyStatement}, \lean{MultVerifyCompletenessStatement} & the check and $\Fn{multVerify}$ & \ref{sec:gadgets}\\
\lean{MainStatement}, \lean{MainCompletenessStatement} & one batch and $\Fn{main}$ & \ref{sec:gadgets}\\
\lean{TprfCorrectnessStatement}, \lean{TprfSecurityStatement} & the threshold PRF & \ref{sec:sign}\\
\lean{SigningCorrectnessStatement} & an honest quorum signs & \ref{sec:sign}\\
\lean{SigningSecurityStatement} & strong unforgeability, adaptive corruption & \ref{sec:sign}\\
\lean{LifetimeSecurityStatement} & the composed lifetime, not in the bundle & \ref{sec:lifetime}\\
\bottomrule
\end{tabular}
\end{center}

\subsection*{The literature}

The DKG follows Goyal and Song (GS20, ePrint 2020/134) and the note \lean{doc/leanXMSS.pdf}. The
threshold PRF and the signing follow PRAWNS (Boneh, B\"unz, D., Litchev and Nazirkhanova, ePrint
2026/2142). The DKG claims use the simulation framework of Canetti (2000) and of Goldwasser and
Lindell (GL05). They have selective abort and static corruption, as GS20 proves. The signing claims
follow PRAWNS, Definitions 2.2 and 3.2, with adaptive corruption, and leanVM for strong
unforgeability. The other keys below are these.
\begin{itemize}\itemsep0pt
\item CDI05 (Cramer, Damg\aa rd and Ishai) gives replicated sharing and the PRSS.
\item CDG+18 gives the local programmable random oracle, and CGH+18 gives the agreement check.
\item KPPS23 is MPClan.
\item BR93 gives the random-oracle model and its query bounds.
\item BDG20 (Bellare, Davis and G\"unther) gives domain separation.
\end{itemize}

\section{The setting}

\subsection*{The committee}

There are $n = 2t - 1$ operators $P_0, \dots, P_{n-1}$ (\lean{Committee}). Any $t$ of them form a
quorum and sign together. At most $t - 1$ are corrupt (\lean{Coalition}). A \emph{term} $T$ is a set
of $t - 1$ operators, and its $t$ \emph{holders} are the operators outside it. There are
$\NT = \binom{n}{t-1}$ terms and $\NP = \binom{n}{2}$ pairs of operators. A coalition of $t - 1$
operators holds every term except one, the term equal to itself. So at least one term stays hidden
from every coalition, and every term has an honest holder. The rounds are synchronous, and the links
are private and authenticated (GS20, Section 3.1).

\subsection*{Replicated sharing and the PRSS}

A wire value $v$ is a replicated sharing $v = c + \sum_T v_T$ (CDI05). The part $c$ is public, and
the holders of $T$ hold $v_T$. A Shamir sharing over the check field is used only inside the check.
Fresh random sharings come from the PRSS of CDI05. Each term key and each pair key of a ceremony
gives a stream $\mathsf{PRG}(\mathrm{key}, \rho, \mathrm{idx})$, so a random sharing costs no
message. The PRSS outputs of one round realize $\Fn{rand}$, which has no abort, because the PRSS
sends no message. Each ceremony runs its own setup for fresh PRSS keys.

\subsection*{leanXMSS and the keys}

The signature is leanXMSS at the parameters of leanVM. It has $58$ Winternitz chains of length $4$,
a $128$-bit message digest and a target sum of $106$. Its two layers of trees of height $h = 16$ give
$2^{32}$ one-time keys. One signature makes at most $L = 2^{32}$ randomizer trials. The keygen
deals three kinds of keys by commit, reveal and echo.
\begin{itemize}\itemsep0pt
\item The public parameter $P$.
\item A seed $\mathrm{seed}_T$ for each term, which the holders of $T$ know.
\item A pair key $K_{ij} = K_{ji}$ for each pair of operators.
\end{itemize}
The threshold PRF is $F(x) = \bigoplus_T H_{16}(\mathrm{seed}_T \Vert x)$, the first $16$ bytes of
$H$. Its inputs $x$ are the chain positions $A_{\mathrm{prf}}(e, l, p)$ and the randomizer values
$A_{\mathrm{rnd}}(e)$ of each key $e$ (\lean{Point}).

\subsection*{The ceremonies}

The trees are built in $1 + 2^{16}$ ceremonies (\lean{ceremony}), one for the top tree and one for
each bottom tree $b$ (\lean{Tree}). A ceremony computes the chains of its keys in MPC from the
PRF values, as boolean circuits of Keccak-f[1600]. A bottom ceremony also opens, for each key $e$,
the masked chain values $\mathrm{cpk}_e$ and the commitment
$c_e = \Th(P, A_{\mathrm{rnd}}(e, 1), F(A_{\mathrm{rnd}}(e)) \Vert 0)$. It also signs its bottom
root with its top key. A ceremony splits its jobs into batches of \lean{batchSize} jobs. Each batch is one
run of GS20, Protocol 19. It runs the AND layers with $\mathsf{Reshare}$ and rerandomizes its
outputs. Then it runs the GS20 check over the field of $2^{192}$ elements and opens its outputs with
$\mathsf{OpenRep}$, as in MPClan. Each operator aborts alone. The tree dealer $\Fn{tree}$
(\lean{Ftree}) computes the tree in the clear with the same circuits on the seeds of the honest
holders. It leaks $P$ and the openings of each batch.

\subsection*{The signing}

The signing needs no MPC (\lean{doc/sign.md}). A quorum $J$ signs a message $m$ at key $e$ in two
rounds. Each signer $i$ sends the part $y_i = \mathsf{Eval}_i(x, m, J)$. This part is the XOR of the
PRF terms that go to $i$ and one mask $H_{16}(K_{ij} \Vert x \Vert m \Vert J)$ for each co-signer
$j$. $\mathsf{Combine}$ XORs the parts, so the masks cancel in pairs and the result is $F(x)$.
\begin{enumerate}\itemsep0pt
\item Round 1 opens the randomizer value $\rho_e = F(A_{\mathrm{rnd}}(e))$. Each signer aborts unless
it matches $c_e$.
\item Each signer finds the first randomizer $R$ from $\Th(P, A_{\mathrm{rnd}}(e), \rho_e \Vert m)$
whose digest is a codeword. This search has no message.
\item The codeword gives a position $p_l < 4$ on each chain $l$. Round 2 reconstructs
$F(A_{\mathrm{prf}}(e, l, p_l))$ from the parts for each chain with $p_l < 3$. The chain value is this
PRF value at position $0$. At positions $1$ and $2$ it is the public entry $\mathrm{cpk}_e[l, p_l]$ XOR
the PRF value. At position $3$ it is the public chain end $\mathrm{cpk}_e[l, 3]$, with no PRF part.
Each signer aborts unless each chain value leads to its chain end.
\end{enumerate}
The output is a plain leanXMSS signature. Honest operators sign at most one message at each key, over
all quorums (\lean{breaksOneMessagePerKey}). The signing game enforces this rule. The PoC does not
enforce it, and production must enforce it.

\subsection*{The oracles}

$H$ is one oracle on the type \lean{Query}. Each constructor of \lean{Query} is its own domain of
$H$, and a repeated query gets the same answer. The domains are these.
\begin{itemize}\itemsep0pt
\item The tweakable hash $\Th$ (\lean{hash}), the PRF terms (\lean{prf}) and the masks (\lean{mask}).
\item The commitments, the key derivations and the echo hashes of a setup.
\item The copy hashes of $\mathsf{OpenRep}$, and the running histories and salted digests of the check.
\item The stream keys, the coin keys and their expansions.
\end{itemize}
This models domain-separated hashing (BDG20). The model has three worlds, and each world answers
$\Th$ in its own way.
\begin{center}\small
\begin{tabular}{lll}
\toprule
World & $\Th$ on the MPC inputs & Used by\\\midrule
\lean{world} & the fixed $\Th$ of \lean{prims}, free & all claims except the two below\\
\lean{signingWorld} & a lazy random oracle & \lean{SigningSecurityStatement}\\
\lean{lifetimeWorld} & the fixed $\Th$ of \lean{prims}, counted & the lifetime claim\\
\bottomrule
\end{tabular}
\end{center}
The MPC inputs (\lean{MpcInput}) are the chain steps and the commitment $c_e$, the only inputs of
$\Th$ that the circuits evaluate. Each world answers every other input of $\Th$ with a lazy random
oracle, the leaves, the nodes and the message digests included. It answers every other domain in the
same way. The fixed function
$\Th$ is a field of \lean{Primitives}, and the claims quantify over every such function.

\subsection*{The circuits, AES-128 and the check field}

The circuits are a parameter (\lean{Circuits}). The ceremony claims quantify over every circuit family,
because $\Fn{tree}$ runs the same circuits. The hypothesis \lean{CircuitsCompute} states that the
circuits compute the jobs of leanXMSS with the fixed $\Th$. Only \lean{DealerTreeStatement},
\lean{TreeSuccessStatement} and the lifetime claim assume it.

$\mathsf{PRG}(\mathrm{key}, \rho, \mathrm{idx})$ is AES-128 under the first $128$ bits of
$H(\mathrm{key})$, as in the implementation. The model takes AES-128 as any PRF scheme with a
$128$-bit key (\lean{Aes}). The scheme \lean{hiddenStreams} has one independent AES key for each
PRSS key that the coalition does not know. The multi-key PRF advantage against it is a term of each
bound that needs it.

The check field $E$ is any field of characteristic $2$ with a value of \lean{Constants}. The field
\lean{encode} of \lean{Constants} gives $E$ exactly $2^{192}$ elements (\lean{Constants.card}), so $E$
is isomorphic to $\mathrm{GF}(2^{192})$. The weights of the bits of a word are any family that is
linearly independent over $\mathbb F_2$, and the implementation uses the powers $x^q$. No claim fixes
the basis or the bit encoding.

\subsection*{The paper model}

The protocol is modeled at the level at which GS20 and PRAWNS state it. Messages carry values, not
bytes, and the model assumes that every encoding of the implementation is injective. Framing, padding
and the byte layout of messages are not modeled. The digests of the check travel in their own round,
and the implementation sends them with the last coin. The session nonces and the commitment of the
PRSS setup in the nonce round are not modeled, and fresh PRSS parts separate the sessions. A run
has a fixed number of steps (its fuel). The model splits each run into batches of one size, and the
implementation sizes chain batches and commitment batches apart.

\subsection*{Secure computation}

An adversary is rushing and malicious. It gets the inputs of the corrupt operators and an auxiliary
input $z$. A distinguisher with oracle access reads the joint output of the adversary and the honest
operators. With selective abort, the adversary names the honest operators that receive each output
(GL05). The claims use three notions.
\begin{itemize}
\item \lean{SecurelyComputes}, in a hybrid model. For each coalition $D$ and adversary
$\mathcal A$ there is a simulator $S$. It is fixed before the inputs, $z$, the distinguisher and the
budget. It sees only the corrupt inputs, $z$ and the leaks of the functionality, and it cannot program
$H$. Take every valid input and every $q$-query adversary and distinguisher. Then the ideal run makes
at most $q + \mathit{overhead}$ counted calls, and the two runs differ by at most $\varepsilon(q)$.
\item \lean{SimulatedBy}, on the real network. The simulator is the named code $\mathit{Sim}(D,
\mathcal A)$. It runs $\mathcal A$ on the counted copy. It emulates the honest operators on the free
copy as the real operators do. It can program $H$ as the local programmable random oracle of CDG+18. The
same budget and distance conditions hold.
\item \lean{Completes}. An honest run, with no corruption, gives the outputs of the functionality, up
to $\varepsilon$ in statistical distance. GS20 states no completeness lemma, and security with abort
alone allows a protocol that always aborts. So each completeness claim is the honest run of its
protocol.
\end{itemize}

\subsection*{The budgets}

The bounds count oracle calls, not time. There are three budgets.
\begin{itemize}
\item In the simulation claims and in \lean{TprfSecurityStatement}, $q$ bounds the counted calls of
the adversary and of the distinguisher (\lean{IsQueryBounded}, BR93). The bound holds for every
answer of the oracles and every state. The honest operators and the functionalities call $H$ for free.
\item In the signing claim, $q$ bounds every call of $H$ and $\mathsf{PRG}$ by every party, cache hits
included (\lean{HasSigningQueryBound}). The key generation, the signers, the forger and the final
verification all count.
\item In the lifetime claim, $Q$ counts in the same way on every path of the lifetime game, for every
answer of the oracles (\lean{HasLifetimeQueryBound}). It does not bound an adversary state that no
answer reaches.
\end{itemize}

\section{The keygen}\label{sec:keygen}

The keygen (\lean{keygen}) runs $3$ rounds (\lean{setupRounds}). In round 1, each operator sends a
commitment $\mathrm{com}_{i \to j}$ that names its sender and its receiver. In round 2 it sends its
reveal. In round 3 it sends the hash of all parts of $P$ and an echo hash. An operator aborts unless each reveal opens its commitment and each
echo equals its own. Each key hashes its tag and the revealed parts. The functionality $\Fn{keygen}$
(\lean{Fkeygen}) is the trusted dealer of PRAWNS $\mathsf{KeyGen}'$ with selective abort. It samples a
uniform $P$, uniform seeds and uniform pair keys. It leaks $P$ and the corrupt keys, and gives the
keys to the honest receivers that the adversary names.

\paragraph{Completeness} (\lean{KeygenCompletenessStatement}). An honest keygen gives the keys of
$\Fn{keygen}$ up to the statistical distance
\[
\boxed{\ (\NT^2 + \NP^2)\cdot 2^{-256}\ }
\]
(\lean{derivedKeyCollision}). A derivation hashes only its tag and its parts. The bad event is two
terms, or two pairs, with equal derivation inputs. Equal inputs force equal keys, while the dealer
samples the keys independently. Distinct inputs give independent uniform keys. These keys can
collide, as the keys of the dealer can.

\paragraph{Security} (\lean{KeygenSecurityStatement}). Against every static coalition of at most
$t - 1$ operators, the keygen on the real network securely computes $\Fn{keygen}$ with selective
abort, with the explicit simulator \lean{keygenSimulator}. The simulator observes every call of the
adversary. After round 1, it extracts the corrupt reveals from the calls with the commitment prefix.
Then it programs $H$ so that $P$ and the corrupt keys are those of $\Fn{keygen}$. Its overhead is
$1 + \NT + n^2$ programming queries (\lean{keygenOverhead}). The distance is at most
\[
\boxed{\ (q + s_K)^2\cdot 2^{-256}, \qquad s_K = n(1 + \NT + n) + 1 + \NT + \NP + n^2\ }
\]
(\lean{keygenHashError}, \lean{keygenSecrets}). The $s_K$ honest points are the parts of each
operator, the derived keys and the inputs of the echo hashes. The birthday bound of the $q$ calls and
these points pays a guessed secret, a collision of two secrets, a failed programming and a failed
extraction.

\section{The ceremonies}\label{sec:tree}

Each ceremony claim fixes a top tree and a tree, and assumes \lean{batchSize} $> 0$. Its inputs are
the keys of one dealing (\lean{DealtInputs}), because the operators build the bottom trees on demand
while they sign. The coalition of a ceremony is the set of operators corrupted before it, and no
operator is corrupted during it. Let $B$ be the number of checked batches of the tree
(\lean{treeBatches}). Let $d$ and $R$ be the largest numbers of coins and of halving steps of the
check of one batch, over all runs (\lean{maxBatch}).

\paragraph{The check error.} One check with $d$ coins and $R$ halving steps misses a wrong tuple with
probability at most
\[
\mathrm{checkErr}(d, R, q) = \frac{d + 2(R + 1)}{|E| - 2} + \frac{(R + 2)\,q}{|E|}
\]
(\lean{checkErrorOf}, GS20, Lemmas 8 to 10). The first term is the Schwartz--Zippel count. The second
term pays a counted call that hits the key input of one of the $R + 2$ coins before the coin opens.
The ceremony pays $B \cdot \mathrm{checkErr}(d, R, q)$ (\lean{treeCheckError}).

\paragraph{Completeness} (\lean{TreeCompletenessStatement}). An honest ceremony gives the output of
$\Fn{tree}$ up to the statistical distance
\[
\boxed{\ B\left(\frac{7}{|E|} + (R + 1)\left(\frac{2}{|E|}\right)^{2^{64}}\right)
+ \frac{(B(R + 2))^2}{2^{257}}\ }
\]
(\lean{treeAbortError}). The first term is a slice coin equal to $0$, or a challenge that finds no
value outside $\{0, 1\}$ within $2^{64}$ tries (\lean{coinFailure}). The second term is a collision of
two of the $B(R + 2)$ coin keys, because a later check reads the cached coins of an earlier batch.

\paragraph{Security} (\lean{TreeSecurityStatement}). Against every static coalition of at most
$t - 1$ operators, each ceremony on the real network securely computes $\Fn{tree}$ with selective
abort, with the explicit simulator \lean{treeSimulator} and no overhead. The simulator runs the
adversary with emulated honest operators. They use $P$, the corrupt seeds for the terms that a corrupt
operator holds, zero seeds elsewhere, and ideal streams under the hidden PRSS keys. The opening of
each batch shifts the first lacked term $T_0$, so that it opens the output of $\Fn{tree}$ (GS20,
Theorem 2, output phase, and MPClan, Appendix B). $\mathsf{Rerand}$ first makes $T_0$ uniform. For an
adversary $\mathcal A$, inputs $x$, auxiliary input $z$ and distinguisher $Z$, the distance is at most
\[
\boxed{\ \begin{aligned}&B\cdot\mathrm{checkErr}(d, R, q) + \frac{(q + s_T)^2}{2^{256}} + \frac{(\NT + \NP)^2}{2^{129}}\\
&+ \mathrm{Adv}^{\mathrm{prf}}_{\mathtt{hiddenStreams}(\mathrm{AES})}\big(\mathtt{prgReduction}(D,
\mathrm{Real}_{\mathcal A, x, z, Z})\big)\end{aligned}\ }
\]
with the honest points (\lean{treeSecrets})
\[
s_T = n(\NT + n) + \NT + \NP + n^2 + (\NT + \NP) + B\,(3\NT + n\NT + R + 2).
\]
\begin{itemize}
\item The first term is the soundness of the checks of all batches.
\item The second term (\lean{treeHashError}) is the birthday bound of the counted calls and the honest
points. The PRSS setup has the parts, the derived keys and the echo inputs. Each PRSS key has a
stream-key input. Each batch has the history, the digest and the copy hash of each term, the copies
that a corrupt owner sends to each honest operator, and the $R + 2$ coin keys.
\item The third term (\lean{streamKeyError}) is a collision of the $128$-bit AES keys of the
$\NT + \NP$ PRSS keys.
\item The last term is the multi-key PRF advantage of the reduction \lean{prgReduction} on the real
run. The reduction answers each honest stream under a hidden key from the PRF oracle (CDI05, the PRSS
of the lacked term). The term is on the real run because \lean{treeSimulator} already uses ideal
streams.
\end{itemize}

\paragraph{The dealer gives the trees} (\lean{DealerTreeStatement}). Assume that the circuits
compute $\Th$ (\lean{CircuitsCompute}) and that \lean{batchSize} $> 0$. Fix a parameter, seeds, a top
tree, a tree and a cache of $H$. Then the dealer of $\Fn{tree}$ and \lean{referenceTree} with the PRF
values of the seeds give the same distribution of the tree and of the next cache. \lean{referenceTree} is the key
generation of PRAWNS, Section 3.1, for the trees of leanXMSS. So $\Fn{tree}$ gives the trees of the
signing claims.

\paragraph{The dealer rarely fails} (\lean{TreeSuccessStatement}). Under the same hypotheses, the
dealer that builds the top tree and then bottom tree $b$ on it fails with probability at most
\[
\boxed{\ (1 - p)^{L}, \qquad L = 2^{32}\ }
\]
(\lean{searchFailure}). Here $p$ (\lean{codewordProbability}) is the probability that one trial of
the randomizer search hits a codeword, the share of $128$-bit digests whose digits sum to $106$. The
search is that of the top signature of the bottom root.

\section{The gadgets}\label{sec:gadgets}

The gadget claims are the lemmas of GS20 for replicated sharing. They are internal steps, and the
keygen and ceremony claims do not depend on their interfaces. Each is in a hybrid model, and the
PRSS realizes $\Fn{rand}$. Their functionalities leak the view of each honest operator of the terms
that the corrupt operators know, more than GS20 on some inputs. Every simulator of this section
exists by the claim, sees only the leaks, and has all its calls counted.

\paragraph{The valid inputs.} Each gadget claim is stated only on valid inputs. Below, $i$ and $j$ are
honest operators, and $x_i$ is the input of $i$.
\begin{itemize}
\item \lean{SameCount} (the coin and the multiplications) states $x_i = x_j$ for the count of values.
\item \lean{ShamirInputs} (\lean{ExtendMultStatement}) states that $i$ and $j$ agree on the kind and the
length of each entry. For each Shamir entry $g$ and position $k$ there are polynomials $V, W$ of
degree at most $t - 1$ with $a_i(k) = V(\alpha_i)$ and $b_i(k) = W(\alpha_i)$ at every honest $i$.
\item \lean{TupleInputs} (the check) states that the honest lists of tuples have the given lengths and
one mode for each tuple. Fix a term $T$. If the honest holders of $T$ agree on the values $z_T(g)$ that
a corrupt operator fills ($\mathrm{filler}(T, g) \in D$), then they agree on all of $x_T$, $y_T$ and
$z_T$. So a term with no corrupt holder always agrees, and a split copy needs a split value that a
corrupt operator filled.
\item \lean{MainInputs} (one batch) states that there are widths $w_k$, one for each job $k$. The view
of $i$ of term $T$ in job $k$ has length $w_k$ if $i$ holds $T$ and length $0$ otherwise. Every holder
of each term agrees on it, the corrupt holders included.
\end{itemize}

\paragraph{The leaks of the check.} $\Fn{multVerify}$ (\lean{FmultVerify}) takes each term of each
tuple from its first honest holder. Let $\mathit{agree}$ be the event that every two honest holders
of each term have the same $(x_T, y_T, z_T)$ for every tuple. Let $\mathit{accept}$ be the event
$\mathit{agree}$ and $z(g) = x(g)\,y(g)$ for every reconstructed tuple and position. The functionality
leaks three things.
\begin{itemize}\itemsep0pt
\item The view of each honest operator of the terms that the corrupt operators know.
\item The differences $z(g) - x(g)\,y(g)$ of the reconstructed tuples, only when $\mathit{agree}$
holds. Reconstructed values of split copies can depend on secrets.
\item The bit $\mathit{accept}$.
\end{itemize}
An honest receiver gets an output only when $\mathit{accept}$ holds.
\begin{center}\footnotesize
\begin{tabular}{llll}
\toprule
Claim & Protocol and model & Overhead & Error\\\midrule
\lean{CoinStatement} & $\mathsf{Coin}$, $\Fn{rand}$-hybrid (GS20, Lemma 3) & $0$ & $0$\\
\lean{MultStatement} & $\mathsf{MulRep}$, $\mathsf{Reshare}$, each mode (Lemma 4) & $0$ & $0$\\
\lean{ExtendMultStatement} & $\mathsf{MulRep}$ or $\mathsf{MulSh}$, $\mathsf{Reshare}$ (Lemma 5) & $0$ & $0$\\
\lean{MultVerifyStatement} & the check, $\{\Fn{coin}, \Fn{extendMult}, \Fn{rand}\}$ (Lemmas 8 to 11) & $c_V$ & $\varepsilon_V(q)$\\
\lean{MainStatement} & one batch, $\{\Fn{mult}, \Fn{multVerify}\}$ (Theorem 2) & $2n\NT$ & $\varepsilon_M(q)$\\
\bottomrule
\end{tabular}
\end{center}
The multiplications are secure up to an additive attack, as in GS20. For the check of $k$ tuples with
$d$ coins and $R$ halving steps,
\[
c_V = n\big(R + 2 + d + (R + 1)\,2^{64} + 2\NT\big), \qquad
\varepsilon_V(q) = \mathrm{checkErr}(d, R, q) + \frac{(q + 2n\NT + R + 2)^2}{2^{256}}
\]
(\lean{checkOverhead}, \lean{checkHashError}). The overhead is the emulation of the coin keys, the
expanded coins, the challenge tries and the two hashes of each term at each operator. The second term
of $\varepsilon_V$ is the birthday bound of the counted calls, the histories and digests of the terms,
and the coin keys. For one batch,
\[
\varepsilon_M(q) = (q + 2n\NT)^2\cdot 2^{-256}
\]
(\lean{copyHashError}). The adversary chooses the honest copy of each term that it holds at the
$\mathsf{Rerand}$ leak, so a collision of two copies passes the hash check of $\mathsf{OpenRep}$.
\lean{MainStatement} is GS20, Theorem 2, for public outputs and for inputs that the operators already
share. Its inputs agree at every holder, as the keygen echo makes them.

\begin{sloppypar}\noindent The completeness claims are the honest runs. \lean{CoinCompletenessStatement} and
\lean{MainCompletenessStatement} hold with distance $0$. \lean{MultVerifyCompletenessStatement}
states that an honest check accepts correct tuples (\lean{CorrectTuples}) up to the distance
$\mathrm{coinFailure}(R) + (R + 2)^2/2^{257}$ (\lean{abortErrorOf}).\end{sloppypar}

\paragraph{The deviations from GS20.} \lean{Statement.lean} lists them, and each has its own argument.
\begin{itemize}\itemsep0pt
\item Replicated sharing replaces Shamir sharing. A corrupt filler can give each honest holder its own
copy, so the salted digests compare the copies (KPPS23, $\Pi_{\mathrm{agree}}$, and CGH+18, Protocol
3.3).
\item Each operator aborts alone, so every functionality except $\Fn{rand}$ has selective abort.
\item The de-linearization weighs tuple $k$ by a product of $d$ coins, as the prototype does, in
place of the powers of one coin of GS20, Protocol 13.
\item Bit $q$ of a word has the public weight $x^q$. This is the partially deterministic challenge of
Dao and Thaler (2024), with the geometric weights of Flock (B\"unz, Rothblum and Wang, 2026).
\item The first seven halving steps fold slices with the coin kept out (Gruen, 2024) and a table of
small values (BDDT25). Each step is still a degree-$2$ claim in its challenge.
\item GS20 assumes exactly $(n-1)/2$ corrupt parties. \lean{treeSimulator} handles every coalition
size, because it shifts one lacked term at the opening.
\end{itemize}

\section{The threshold PRF and the signing}\label{sec:sign}

Only \lean{SigningSecurityStatement} runs in \lean{signingWorld}, where all of $H$ is a lazy random
oracle, $\Th$ included. The two PRF claims and \lean{SigningCorrectnessStatement} run in
\lean{world}. There the PRF terms and the masks are random-oracle domains, and $\Th$ is fixed on the
MPC inputs. The committee is a parameter of each claim, and PRAWNS lets the adversary choose it.
$\mathsf{Eval}'$ of PRAWNS, Section 5.1, adds two masks under $k_{ij}$ and $k_{ji}$. The note adds one
mask under the common key $K_{ij}$, and the masks still cancel in pairs.

\paragraph{PRF correctness} (\lean{TprfCorrectnessStatement}). For every quorum $J$, input $x$ and
context $m$, $\mathsf{Combine}'$ of the parts of $J$ gives $F(x)$ with probability $1$ (PRAWNS,
Definition 2.1).

\paragraph{PRF security} (\lean{TprfSecurityStatement}). The game $G^{\mathrm{tprf}}$ of PRAWNS,
Figure 1 (\lean{tprfGame}), deals the keys with $\mathsf{KeyGen}'$ of PRAWNS, Section 5.1. The
adversary has the oracles $\mathsf{OCorrupt}$, $\mathsf{OEval}$ with the quorum and the context, and
$\mathsf{OChallenge}$. A challenge on $x$ gives $F(x)$ or a random value. The game refuses a query
that brings the corrupt operators and the parts of one $(x, m)$ to $t$ on a challenged point. For
every adversary with at most $q$ counted calls,
\[
\boxed{\ \mathrm{Adv}^{\mathrm{tprf}} \le \frac{(2\NT + n^2)\,q + (\NT + n^2)^2}{2^{256}}\ }
\]
(\lean{tprfError}). This is PRAWNS, Theorem 5.1, with $F(k, x) = H_{16}(k \Vert x)$ in the
random-oracle model. PRAWNS bounds it by $2\NT\,\varepsilon_F$. The masks add $n^2 q/2^{256}$ for the
guess of a hidden pair key. Honest calls are free, so two of the $\NT + n^2$ keys can collide.

\paragraph{Signing correctness} (\lean{SigningCorrectnessStatement}). For every quorum, key and
message, every signature that an honest quorum outputs verifies. The quorum fails with probability at
most
\[
\boxed{\ 2\,(1 - p)^{L}\ }
\]
for the search of the top signature in $\mathsf{KeyGen}$ and the search of the signature.

\paragraph{The signing game} (\lean{tsigGame}). The game $G^{\mathrm{tsig}}$ of PRAWNS, Figure 3,
deals the keys with the dealer $\mathsf{KeyGen}$ (\lean{dealerKeyGen}) and gives the forger the
public key and the top tree. The forger has $H$ and the oracles \lean{tsigOracles}.
\begin{itemize}\itemsep0pt
\item $\mathsf{OCorrupt}$ is adaptive and stops at $t - 1$ operators. A corruption gives only the keys
of an operator, because an operator erases its ceremony state.
\item The tree oracle gives the public outputs of a bottom tree and the opened randomizer value of its
top key. The game builds each bottom tree once.
\item $\mathsf{OSign}$ has two rounds and takes the quorum. An honest operator answers only for one
message at each key (\lean{breaksOneMessagePerKey}). This rule replaces line 8 of PRAWNS, Figure 3,
which fails for $n = 2t - 1$.
\end{itemize}
A forgery wins (\lean{forgeryWins}) if it verifies and is fresh. It is fresh in two cases. In the
first case, fewer than $t$ operators can give chain parts at its key for its message. In the second
case, it differs from the signature of $\mathsf{KeyGen}$ at its key. The second case is strong unforgeability, which goes beyond PRAWNS. The
commitment $c_e$ fixes the randomizer value, so the same message at the same key gives the same
signature.

\paragraph{Signing security} (\lean{SigningSecurityStatement}). For every forger and every $q$ that
bounds every call of $H$ and $\mathsf{PRG}$ in the game,
\[
\boxed{\ \mathrm{Adv}^{\mathrm{tsig}} \le \frac{q}{2^{127}} + \frac{(\NT + n^2)\,q}{2^{256}} +
\frac{(\NT + n^2)^2}{2^{256}}\ }
\]
(\lean{signingHashError}, \lean{seedHitError}, \lean{seedCollisionError}).
\begin{itemize}
\item The first term gives each call two units of $2^{-128}$. One unit pays a call whose output hits an
honest output at its tweak. The other pays a call that meets a hidden honest input. Such an input is a
chain value below an opened position, $\rho_e$ before round 1, or the output of a digest that a later
search selects.
\item At most $t - 1$ corrupt operators leave the seed of at least one term and the pair keys of the
honest pairs hidden. The second term pays a call on such a key. Each call names one key, which is
uniform on $256$ bits.
\item The third term pays two equal keys. Equal term seeds cancel in $F$, and a hidden key equal to a
corrupt key is known.
\end{itemize}
Without these events, each unopened $F(x)$ is a fresh uniform value. This covers $\rho_e$ before
round 1, the chain starts and the masks.

\paragraph{What the signing claims do not give.} The claim is in the random-oracle model for all of
$H$, $\Th$ included, while the ceremonies evaluate a fixed $\Th$ in MPC. The step from one model to
the other is the random-oracle heuristic for Keccak. The claim also takes the trees of the dealer
$\mathsf{KeyGen}$, not the trees of the real ceremonies. Section~\ref{sec:lifetime} composes the two.
The one-message rule is a named coordination assumption, and the PoC does not meet it.

\section{The lifetime claim}\label{sec:lifetime}

\lean{LifetimeSecurityStatement} states the security of the real keygen, the real ceremonies with
retries and the signing service together. It fixes one $\Th$ (the field of \lean{prims}) and one
circuit family, and it assumes \lean{CircuitsCompute} and \lean{batchSize} $> 0$.

\paragraph{The lifetime game} (\lean{lifetimeGame}). The coalition $D$ of at most $t - 1$ operators is
static and is corrupt from the start.
\begin{itemize}
\item The keygen runs once on the real network. If an honest operator has no keys, the lifetime ends.
The adversary then gets the public parameter.
\item The adversary gets the public key and the top tree of the first honest operator $h$ once $h$
holds its top tree.
\item The adversary runs one ceremony at a time and chooses each tree and the start of each run. Each
tree has at most \lean{attempts} runs. A run of a tree starts only while some honest operator lacks
that tree. A retry runs on the same keys with a fresh PRSS setup. Each operator keeps its first output
of each tree, and a later run does not replace it.
\item A bottom ceremony runs only when every honest operator has the top tree, and not after the first
signing request on its tree.
\item Signing requests can come between the rounds of a ceremony. A request on the tree of the running
ceremony is refused. An honest operator answers a request, round 1 included, only when it holds its own
ceremony output of that tree. It signs at most one message at each key.
\item The adversary loses if an honest operator has no top tree at the end. The win test is
\lean{forgeryWins}, with $\mathsf{KeyGen}$ on the seeds of the honest operators.
\end{itemize}
The game runs in \lean{lifetimeWorld}. $H$ answers the MPC inputs with the fixed $\Th$ and every other
input as a lazy random oracle.

\paragraph{The bound.} Let $Q$ bound every call of $H$ and $\mathsf{PRG}$ on every path of the game.
Let $q = (\mathit{fuel} + 2)\,Q$ with the fuel of a ceremony (\lean{ceremonyFuel}),
$q_0 = Q + \mathit{ov}(Q)$, and
\[
\mathrm{run}(\mathit{tree}) = \mathit{attempts}\cdot\Big(B\cdot\mathrm{checkErr}(d, R, q_0 + q) +
\frac{(q_0 + q + s_T)^2}{2^{256}} + \frac{(\NT + \NP)^2}{2^{129}} +
\frac{(\NT + n^2)(q_0 + q)}{2^{256}}\Big)
\]
with $B$ and $s_T$ of the tree. Then
\[
\boxed{\ \begin{aligned}
\Pr[\text{the lifetime adversary forges}] \le{}& \frac{(5Q + s_K)^2}{2^{256}}
+ \mathrm{run}(\mathrm{top}) + \sum_{b < 2^{16}} \mathrm{run}(b)\\
&+ \sum_{r < N_{\mathrm{run}}} \mathrm{Adv}^{\mathrm{prf}}_{\mathtt{hiddenStreams}(\mathrm{AES})}
\big(\mathtt{lifetimeRunReduction}(r)\big)\\
&+ 2^{16}\,(1 - p)^{L} + \mathtt{lifetimeForgerAdvantage}
\end{aligned}\ }
\]
with $N_{\mathrm{run}} = \mathit{attempts}\,(2^{16} + 1)$ runs (\lean{lifetimeAdvantage}, \lean{lifetimeRuns},
\lean{lifetimeOverhead}).
\begin{itemize}
\item The first term is \lean{keygenHashError} at the budget $(\mathit{setupRounds} + 2)\,Q = 5Q$. One step
of a phase adversary holds all that the lifetime adversary does between two rounds. So a phase of
$f$ rounds takes the budget $(f + 2)\,Q$.
\item Each run of each ceremony pays its check, its birthday bound and its stream keys, after the
$q_0$ calls of the earlier phases. It also pays \lean{seedHitError} for a call on a hidden key of an
honest term or pair. The emulated signing keeps the masks of the honest pairs in its own table.
\item Each run pays the PRF advantage of its own reduction \lean{lifetimeRunReduction}. That reduction
runs \lean{prgReduction} on a hybrid in which the earlier runs use ideal streams and the later runs
use real streams. The keygen and the signing call no $\mathsf{PRG}$.
\item Each bottom tree pays one failed search, because each retry gives the same search.
\item The overhead $\mathit{ov}(Q)$ is the cost of the trees that a reduction computes in the clear.
\[
\mathit{ov}(Q) = \big(\min(Q, 2^{16}) + N_{\mathrm{run}}\big)\,b + N_{\mathrm{run}}\,n\,b + 1 + \NT + n^2
\]
Here $b$ (\lean{bottomTreeCalls}) bounds the calls of \lean{referenceTree} for one tree, with $\NT$
calls for each PRF value.
\[
b = 2^{16}\cdot 58\,(3\NT + 6) + 2^{16}(\NT + 1) + 58\,(\NT + 6) + \NT + 2^{18} + 2^{32}
\]
Its terms pay the chain rows, the commitments, the Merkle tree, at most $2^{32}$ randomizer trials,
the top signature and its check. The
signing oracles build a tree only after a ceremony run, and each such tree costs at least one call of
$Q$. Each run pays one more tree. The emulated \lean{treeSimulator} of each run evaluates the
openings at each of the $n$ operators. \lean{keygenSimulator} adds its programming queries.
\end{itemize}

\paragraph{The last term.} \lean{lifetimeForgerAdvantage} is the probability that the fixed forger
\lean{lifetimeForger} wins the signing game \lean{tsigGame}, run in \lean{lifetimeWorld} and not in
\lean{signingWorld}. That forger runs the lifetime adversary in a lifetime where the keygen is
$\Fn{keygen}$ and each ceremony is $\Fn{tree}$.
\begin{itemize}\itemsep0pt
\item It corrupts the coalition $D$ and asks no other corruption.
\item It emulates the keygen with \lean{keygenSimulator} on the public parameter of the game. It keeps
the programmed keygen points in its own table, because it cannot program $H$.
\item It emulates each ceremony with the operators of \lean{treeSimulator} and the fixed circuits. It
takes the output of $\Fn{tree}$ from the top tree of the game and from its tree oracle.
\item It gets the keys of the coalition, its pair keys with honest operators included. It gets no
hidden term seed and no key that only honest operators share.
\item It sends the admitted signing requests to the signing oracles of the game, and outputs the
forgery of the lifetime adversary.
\end{itemize}
The forger is fixed code, so the prover does not choose this term. The term has no budget, and no
claim bounds it. \lean{SigningSecurityStatement} does not bound it, because that claim is stated in
\lean{signingWorld}. There all of $H$ is a random oracle, while this forger runs circuits of the fixed
$\Th$. The step from this term to a concrete bound is the random-oracle heuristic for $\Th$. Under
that heuristic, \lean{SigningSecurityStatement} applies to \lean{lifetimeForger} at its own number of
calls. No claim states that number or this step. The heuristic is not a term of the bound.

\paragraph{What the lifetime claim does not give.} The coalition is static. An adaptive lifetime claim
does not follow from these claims, although the threshold PRF and signing claims are adaptive. Ceremonies that run
at the same time as another ceremony are outside the game. The claim does not bound its last term.

\section{What is not proved}

\begin{itemize}
\item The claims are stated for the paper model, not for bytes. Messages carry values, and every byte
encoding is assumed injective. The model has no framing, padding, byte layout or tags on the wire. It
uses one batch size, and the implementation sizes chain batches and commitment batches apart. The
session nonces and the commitment of the nonce round are not modeled. Each run has a fixed fuel. The
implementation sends the digests with the last coin, and the model sends them in their own round.
\item The claims assume private authenticated channels ($\Fn{auth}$). The PoC uses plain TCP with no
transport authentication.
\item The model separates the constructors of \lean{Query}, and the implementation does not yet do
so. The threshold PRF and the masks hash a key first with no tag. To match the model, the
implementation needs injective encodings with a distinct tag for each domain and fields of fixed
width or with a length prefix.
\item The keygen, each ceremony and the lifetime claim have a static coalition. Only the threshold PRF
and signing claims cover adaptive corruption, and no adaptive lifetime claim follows.
\item No claim gives liveness or resistance to denial of service. Security with selective abort
allows a protocol that aborts, and the completeness claims cover only honest runs. Frame-length
attacks are out of scope.
\item The random-oracle heuristic for $\Th$ and for Keccak is not proved. The signing claim takes all
of $H$ as a random oracle. The ceremony and lifetime claims take a fixed $\Th$ on the MPC inputs and a
random oracle elsewhere. No claim treats Keccak-256 or Keccak-f[1600] as the deployed functions.
\item No claim bounds \lean{lifetimeForgerAdvantage}. The lifetime bound holds the advantage of
\lean{lifetimeForger} as an explicit term with no budget.
\item The one-message rule is an assumption on coordination. The signing and lifetime games enforce
\lean{breaksOneMessagePerKey}. The PoC does not enforce it, and production must.
\item \lean{CircuitsCompute} is a hypothesis. No claim shows it for the circuits of the
implementation.
\item AES-128 is not shown secure. The model takes any PRF scheme, and its multi-key PRF advantage is
a term of the bounds.
\item No claim fixes the constants of the implementation. The claims take every value of
\lean{Constants}, and \lean{constants_nonempty} gives a value only for the committee of one operator.
\item Each operator is assumed to erase its ceremony state at the end of each ceremony, also on
abort. Only this erasure makes a corruption give only the keys.
\item A ceremony that runs during another ceremony is outside the claims. In the separate ceremony
claims, a ceremony during an open signing session is outside them too.
\item The gadget functionalities leak more than GS20 on some inputs, and their claims restrict the
inputs. These claims are internal steps.
\item The budgets count oracle calls, not running time. The adversary is classical, and no claim
covers quantum adversaries.
\item No test links the Rust implementation to \lean{Scheme.lean}.
\end{itemize}

\end{document}
